\section{Separate bit fields from Huffman symbols}
A bit reservoir holds unread compressed bits in a register. Track its count separately
from the input cursor. To obtain an ordinary $n$-bit field, refill with bounded byte
loads until at least $n$ bits are available, take the low $n$ bits and shift the
reservoir right by $n$. The block header has a one-bit final flag and a two-bit type.
Type zero is stored, one uses fixed Huffman codes, two supplies code lengths, and
three is invalid. A nonfinal block continues into the next block header.

Huffman symbols need a different interpretation of that arriving bit sequence.
Code bits arrive in code order; do not read an entire code through the ordinary
low-bit numerical helper and compare it directly with an unreversed canonical code.
For a simple illustrative alphabet with codes \texttt{0}, \texttt{10}, \texttt{11},
the sequence \texttt{1,0} selects the second symbol. The ordinary two-bit numerical
interpretation of the same arrivals is one, not two. Keep bit arrival order and code
numbering separate when writing the symbol decoder.

A stored block discards the remaining bits of the current byte before reading its
little-endian length and complement. A compressed block does not discard bits between
symbols. Final block completion also does not authorize another input byte: account
for unused bits already fetched, and reject extra compressed bytes under the exact
stream contract. A refill strategy must not blur bytes in the reservoir with bytes
not yet consumed from the input window.
